% % % % % % % % % % % % % % % % % % % % % % % % % % % % % % % % % % % % % % % % % % % % % % % % % % % % % % % % % % % % % % % %
% SAT_STYLE_MATH_MCQS.tex
% 1_10.tex
% 
% encoding: UTF-8 
% EOL: LF
%
% format: LaTeX
% indent: spaces (2)
% width: 127
% % % % % % % % % % % % % % % % % % % % % % % % % % % % % % % % % % % % % % % % % % % % % % % % % % % % % % % % % % % % % % % %
Considérons deux entiers $a$ et $b$ dont le produit vaut $55$. Sachant que $a-b = \minus 6$, que vaut le doublet $(a, b)$?
%
\begin{enumerate} 
  \item $(1,55)$
  \item $(11,5)$
  \item $(\minus 11,5)$
  \item $(11, \minus 5)$
  \item $(\minus 11,\minus 5)$
  \item $(\minus 5,\minus 11)$
  \item Autre réponse
\end{enumerate}
% 
\begin{proof}%
Deux approches - arithmétique et algébrique - sont valables. La première consiste en une étude de cas, la seconde ramène le %
problème à la recherche de racines d'un polynôme particulier. 
%
\paragraph{Approche arithmétique.} Notons que $m= \min\bigl\{\abs{a}, \abs{b}\bigr\} > 1$ (en effet, $m = 0 \implies ab=0$, %
et $m = 1 \implies \abs{a - b} = 54$). La décomposition de $55$ en facteurs premiers est $(5, 11)$. %
Les arrangements possibles pour $\magnitude*{a}$ et $\magnitude*{b}$ sont donc soit $(5, 11)$, soit $(11, 5)$. %
Comme $\sgn(a) = sgn(b)$ (car $ab > 0$) et que $a - b < 0$, les options restantes sont:
%
\begin{itemize}[label=$\circ$]
  \item $(a, b) = (5, 11)$ \qquad (convient, car $5-11 = \minus 6$),
  \item $(a, b) = (\minus 11, \minus 5)$ \qquad (convient, car $\minus 11 - (\minus 5) = \minus 6$).
\end{itemize} 
%
Le problème a donc deux solutions: réponse \textbf{G}.
%
\paragraph{Approche algébrique} Recherchons $a$ et $\minus b$ vérifiant
%
\begin{equation}
  \begin{cases}
    a \cdot (\minus b) = \minus 55 \\
    a + (\minus b) = \minus 6.
  \end{cases}
\end{equation}
De façon équivalente, $\{a, \minus b\}$ forme l'ensemble des racines de $P(X) = X^2 + 6 X - 55$. Le discriminant de $P(X)$ est 
%
\begin{align}
  \Delta &= 6^2 - 4 \cdot 1 \cdot (\minus 55) &&\\
  &= 36 + 220  && \\
  &= 256  && \\
  &= 16^2    &[\text{exactement deux racines}]
\end{align}
%
Les couples $(a, b)$ possibles sont donc $(\minus 11, \minus 5)$ et $(5, 11)$. Réponse \textbf{G}.
\end{proof}
%